\documentclass[12pt,a4paper]{article}
\usepackage[utf8]{inputenc}
\usepackage[T1]{fontenc}
\usepackage[italian]{babel}
\usepackage{libertinus}
\usepackage{amssymb}
\usepackage[a4paper,top=25mm,bottom=25mm,left=22mm,right=22mm,headheight=15pt]{geometry}
\usepackage{graphicx}
\graphicspath{{../../assets/}{../../Assets/}{assets/}{Assets/}}
\usepackage[table]{xcolor}
\usepackage{tabularx,array}
\usepackage{titlesec}
\usepackage{titletoc}
\usepackage{microtype}
\usepackage{fancyhdr}
\usepackage{enumitem}
\usepackage[hidelinks]{hyperref}

% Numero del verbale: da impostare una sola volta (usato in intestazione, copertina e ID)
\newcommand{\vnum}{7}

% Palette colori del gruppo Bending Forks
\definecolor{bfblue}{HTML}{1B3F9E}
\definecolor{bfrow}{HTML}{F2F5FC}
\definecolor{bfline}{HTML}{C9D1E6}
\definecolor{bfhead}{HTML}{2E6BE6}

\newcommand{\tightlist}{\setlength{\itemsep}{1pt}\setlength{\parskip}{0pt}}
\setlist{topsep=2pt}
\setlist[itemize]{leftmargin=1.3em,label=\textbullet}
\setlist[enumerate]{leftmargin=1.6em}
\setlength{\parindent}{0pt}
\sloppy
\setlength{\parskip}{5pt}
\renewcommand{\arraystretch}{1.45}
\setcounter{secnumdepth}{3}
\setcounter{tocdepth}{3}

% Formattazione titoli sezioni
\titleformat{\section}{\LARGE\bfseries\color{bfblue}}{\thesection.}{0.3em}{}[{\color{bfblue}\vspace{-2pt}\titlerule[0.9pt]}]
\titleformat{\subsection}{\Large\bfseries\color{bfblue}}{\thesubsection.}{0.3em}{}
\titleformat{\subsubsection}{\normalsize\bfseries}{\thesubsubsection.}{0.3em}{}
\titlespacing*{\section}{0pt}{14pt}{6pt}
\titlespacing*{\subsection}{0pt}{10pt}{4pt}
\titlespacing*{\subsubsection}{0pt}{10pt}{3pt}

% Indice cliccabile con puntini/trattini
\contentsmargin{2.2em}
\newcommand{\bfleader}{\titlerule*[0.32em]{\textendash}}
\titlecontents{section}[0pt]{\addvspace{5pt}\large\bfseries}{\thecontentslabel.\ }{}{\hspace{0.4em}\bfleader\contentspage}
\titlecontents{subsection}[1.2em]{\addvspace{0pt}}{\thecontentslabel.\ }{}{\hspace{0.4em}\bfleader\contentspage}
\titlecontents{subsubsection}[2.4em]{\addvspace{0pt}}{\thecontentslabel.\ }{}{\hspace{0.4em}\bfleader\contentspage}

% Intestazioni e piè di pagina
\pagestyle{fancy}
\fancyhf{}
\fancyhead[L]{\footnotesize\sffamily\bfseries\color{bfhead}\textls[120]{BENDING FORKS}}
\fancyhead[R]{\footnotesize\sffamily Verbale Interno N. \vnum}
\fancyfoot[C]{\small Pagina \thepage}
\renewcommand{\headrulewidth}{0.9pt}

% Colorazione alternata per le tabelle
\newcommand{\bftable}{\rowcolors{2}{bfrow}{white}}
\usepackage{etoolbox}
\AtBeginEnvironment{tabularx}{\bftable}

\begin{document}

% ---------- Copertina ----------
\begin{titlepage}
\centering
\vspace*{4mm}
\includegraphics[width=105mm]{logo-bending-forks-trasparente.png}\\[4mm]
\underline{bending.forks.swe@gmail.com}

\vspace{42mm}
{\fontsize{24}{28}\selectfont\bfseries Verbale Interno N. \vnum\par}
\vspace{4mm}
{\Large\color{bfblue} Riunione del 12 novembre 2026\par}

\vfill
\begin{tabular}[t]{r|l}
\textbf{Responsabile} & Nicholas Salmaso \\
\textbf{Segretario} & Filippo Rampazzo \\
\textbf{Moderatore} & Filippo Rampazzo \\
\textbf{Verificatore} & Matteo Wu \\
\end{tabular}
\hspace{14mm}
\begin{tabular}[t]{r|l}
\textbf{Versione} & 1.0 \\
\textbf{Tipo} & Verbale Interno \\
\textbf{Destinatari} & \emph{Gruppo Bending Forks} \\
 & \emph{Prof. Tullio Vardanega} \\
 & \emph{Prof. Riccardo Cardin} \\
\end{tabular}

\vspace{12mm}
{\large\emph{\textbf{Abstract dei contenuti:}}}\\[1mm]
\begin{minipage}{150mm}\centering Sintesi e decisioni operative emerse durante la riunione interna del team del giorno 12 novembre 2026, riguardante la pianificazione delle attività e l'avanzamento del progetto.\end{minipage}
\vspace*{6mm}
\end{titlepage}

% ---------- Indice ----------
\vspace*{2mm}
{\fontsize{22}{26}\selectfont\bfseries Indice\par}
\vspace{6mm}
\makeatletter\@starttoc{toc}\makeatother
\newpage

% ---------- Registro delle versioni ----------
\hypertarget{registro-delle-versioni}{%
\section{Registro delle versioni}}

\begin{center}
\begin{tabularx}{\linewidth}{
  >{\hsize=0.55\hsize\raggedright\arraybackslash}X
  >{\hsize=0.75\hsize\raggedright\arraybackslash}X
  >{\hsize=0.85\hsize\raggedright\arraybackslash}X
  >{\hsize=0.85\hsize\raggedright\arraybackslash}X
  >{\hsize=2.00\hsize\raggedright\arraybackslash}X
}
\rowcolor{bfblue}\textcolor{white}{\textbf{Versione}} & 
\textcolor{white}{\textbf{Data}} & 
\textcolor{white}{\textbf{Autore}} & 
\textcolor{white}{\textbf{Verificatore}} & 
\textcolor{white}{\textbf{Descrizione}} \\
1.0 & 2026-11-12 & Filippo Rampazzo & Matteo Wu & Creazione e redazione iniziale del documento \\ \arrayrulecolor{bfline}\hline
\end{tabularx}
\end{center}

% ---------- Informazioni generali ----------
\hypertarget{informazioni-generali}{%
\section{Informazioni generali}}

\hypertarget{durata-e-luogo}{%
\subsection{Durata e luogo}}

\begin{itemize}
\tightlist
\item \textbf{Data:} 12 novembre 2026
\item \textbf{Inizio:} 15:00
\item \textbf{Fine:} 16:55
\item \textbf{Luogo / modalità:} in remoto, tramite Discord
\end{itemize}

\hypertarget{partecipanti}{%
\subsection{Partecipanti}}

\begin{center}
\begin{tabularx}{\linewidth}{>{\raggedright\arraybackslash}X>{\centering\arraybackslash}X>{\centering\arraybackslash}p{2.2cm}>{\centering\arraybackslash}p{2.2cm}}
\rowcolor{bfblue}\textcolor{white}{\textbf{Nome e Cognome}} & \textcolor{white}{\textbf{Ruolo Progetto}} & \textcolor{white}{\textbf{Presente}} & \textcolor{white}{\textbf{Assente}} \\
Marco Beltrame & Amministratore & $\blacksquare$ & $\square$ \\ \arrayrulecolor{bfline}\hline
Lorenzo Milanese & Amministratore & $\blacksquare$ & $\square$ \\ \arrayrulecolor{bfline}\hline
Filippo Rampazzo & Analista & $\blacksquare$ & $\square$ \\ \arrayrulecolor{bfline}\hline
Nicholas Salmaso & Responsabile & $\square$ & $\blacksquare$ \\ \arrayrulecolor{bfline}\hline
Luca Stangherlin & Analista & $\blacksquare$ & $\square$ \\ \arrayrulecolor{bfline}\hline
Aldo Toffoli & Analista & $\blacksquare$ & $\square$ \\ \arrayrulecolor{bfline}\hline
Matteo Wu & Verificatore & $\square$ & $\blacksquare$ \\ \arrayrulecolor{bfline}\hline
\end{tabularx}
\end{center}

% ---------- Ordine del giorno ----------
\hypertarget{ordine-del-giorno}{%
\section{Ordine del giorno}}

\begin{enumerate}
\tightlist
\item \hyperlink{odg-1}{Allineamento preliminare in vista della riunione con la proponente del 13 novembre};
\item \hyperlink{varie-ed-eventuali}{Varie ed eventuali}.
\end{enumerate}

% ---------- Discussione e decisioni ----------

\newpage
\hypertarget{discussione-e-decisioni}{%
\section{Discussione e decisioni}}

\hypertarget{odg-1}{%
\subsection{Allineamento preliminare in vista della riunione con la proponente del 13 novembre}}

\subsubsection{Aggiornamenti sulle Norme di Progetto}
\textbf{Discussione:} \\
Beltrame e Milanese aggiornano il gruppo sulle Norme di Progetto, concentrandosi sulla struttura dei \textit{commit}, sulle norme generali delle \textit{repository} e sulle norme per il versionamento; chiedono inoltre riscontro al gruppo su alcuni dettagli rimasti in sospeso, in particolare sulle sigle di versionamento. Su richiesta di Milanese il gruppo discute e approva regole semplificate per la gestione di piccole modifiche ai documenti concordate a voce. Tutte le decisioni prese in considerazione sono approvate all'unanimità dai presenti.

\textbf{Decisioni prese:}
\begin{itemize}
\tightlist
\item[$\checkmark$] \textbf{[VI-7-01]} Codici di versionamento in versione \textit{x.y.z (major, minor, patch)} per tutti i documenti ad eccezione dei verbali;
\item[$\checkmark$] \textbf{[VI-7-02]} Rimuovere le diciture "Responsabile", "Redattore" e "Verificatore" dalla prima pagina dei documenti che si prevede siano redatti e modificati da più persone nel corso del progetto;
\item[$\checkmark$] \textbf{[VI-7-03]} Le piccole modifiche o correzioni ai documenti concordate a voce possono essere \textit{pushate} direttamente sul \textit{main branch} in deroga alla procedura ordinaria.
\end{itemize}

\textbf{Responsabili attuazione:} Marco Beltrame, Lorenzo Milanese (inserimento in NdP); Amministratori, Verificatori (in corso di progetto)

\textbf{Scadenza:} 2026-11-16

\subsubsection{Aggiornamenti sull'Analisi dei requisiti}
\textbf{Discussione:} \\
Rampazzo mostra al gruppo i requisiti lato software: l'architettura generale del server, il flusso dalla ricezione dell'audio dal dispositivo fino all'LLM, il ciclo ReAct e il ruolo del protocollo MCP; elenca inoltre le tecnologie che potrebbero essere coinvolte nella realizzazione del server.
Toffoli spiega i requisiti lato hardware: selezione di tecnologie, gestione del barge-in, Silero VAD vs. openWakeWord e vincoli computazionali dei dispositivi. Durante la discussione vengono elencate le domande che gli Analisti hanno ritenuto di porre alla proponente. Beltrame e Milanese portano all'attenzione degli Analisti dubbi sulla \textit{context window} e sulla \textit{knowledge base} dei modelli; il gruppo concorda di esporre questi dubbi alla propontente. \\
Il gruppo conduce infine un approfondimento collettivo su \textit{Willow}, indicato dalla proponente come progetto di riferimento.

\subsubsection{Aggiornamenti sul Glossario}
A causa dell'assenza degli incaricati della redazione del Glossario, l'argomento è rimandato alla prossima riunione in cui almeno uno dei due possa riferire in merito.

\newpage

\hypertarget{varie-ed-eventuali}{%
\subsection{Varie ed eventuali}}
Il gruppo concorda di rendicontare le ore di lavoro per ruolo in un file Google Sheets condiviso. \\
Il gruppo fissa la prossima riunione per il 13 novembre 2026, appena dopo il termine della riunione con la proponente, per preparare il diario di bordo del 16 novembre e rinnovare i ruoli di progetto.

\textbf{Decisioni prese:}
\begin{itemize}
\tightlist
\item[$\checkmark$] \textbf{[VI-7-04]} Creare il documento per la rendicontazione delle ore;
\end{itemize}

\textbf{Motivazione:} \\
È necessario ai fini di un'equa assegnazione dei ruoli in futuro.

\textbf{Responsabili attuazione:} Amministratori \\
\textbf{Scadenza:} 2026-11-16

% ---------- Azioni da intraprendere (TODO) ----------
\hypertarget{azioni-da-intraprendere}{
\section{Azioni da intraprendere (TODO)}}

\begin{center}
\begin{tabularx}{\linewidth}{>{\hsize=2.0\hsize\raggedright\arraybackslash}X >{\hsize=1.0\hsize\raggedright\arraybackslash}X >{\hsize=0.5\hsize\centering\arraybackslash}X >{\hsize=0.5\hsize\centering\arraybackslash}X}
\rowcolor{bfblue}\textcolor{white}{\textbf{Descrizione}} & \textcolor{white}{\textbf{Responsabile}} & \textcolor{white}{\textbf{Scadenza}} & \textcolor{white}{\textbf{Priorità}} \\
Aggiornare le Norme di Progetto & Beltrame, Milanese & 2026-11-16 & Alta \\ \arrayrulecolor{bfline}\hline
Creare il documento per la rendicontazione delle ore & Amministratori & 2026-11-16 & Media \\ \arrayrulecolor{bfline}\hline
\end{tabularx}
\end{center}

% ---------- Prossima riunione ----------
\hypertarget{prossima-riunione}{%
\section{Pianificazione prossima riunione}}

\begin{itemize}
\tightlist
\item \textbf{Data proposta:} 13 novembre 2026
\item \textbf{Ora prevista:} 16:00 o appena terminata la riunione con la proponente
\item \textbf{Luogo / modalità:} in remoto, tramite Discord 
\item \textbf{OdG preliminare:} Rinnovo dei ruoli di progetto e preparazione del diario di bordo del 16 novembre 2026.
\end{itemize}

\vspace{35pt}
\noindent
\begin{minipage}[t]{0.45\textwidth}
\textbf{Firma Segretario:}\\[18pt]
\makebox[0pt][l]{\raisebox{0mm}{\includegraphics[width=40mm]{firma_filippo.jpeg}}}\rule{55mm}{0.5pt}\\[2pt]
\small Filippo Rampazzo
\end{minipage}
\hfill
\begin{minipage}[t]{0.45\textwidth}
\begin{flushright}
\textbf{Firma Verificatore:}\\[18pt]
\makebox[0pt][l]{\raisebox{0mm}{\includegraphics[width=40mm]{firma_matteo.jpg}}}\rule{55mm}{0.5pt}\\[2pt]
\small Matteo Wu
\end{flushright}
\end{minipage}

\end{document}